\chapter{Introdução}
\label{chap:intro}

A Internet das Coisas (IoT, do inglês \textit{Internet of Things}) tem se destacado como um campo em expansão, impulsionado pelo uso crescente de ferramentas de automação no cotidiano. Essas ferramentas demandam o controle e a gerência ágeis de dados, tarefas frequentemente vinculadas à leitura de sensores e à atuação em tempo real sobre dispositivos de uso diário. Esse cenário abrange diversas áreas, como segurança patrimonial, controle de veículos autônomos e medicina.

Nesse contexto, a comunicação entre os dispositivos deve atender aos requisitos de dependabilidade e segurança descritos por \citet{Avizienis2004}, pois falhas, indisponibilidade ou comprometimento dos dados podem afetar diretamente o funcionamento dos sistemas e, em aplicações críticas, gerar consequências graves. Além disso, o volume de dados potencialmente sensíveis transmitidos por esses dispositivos torna essencial a confidencialidade das informações. Nesse sentido, mecanismos criptográficos devem proteger o conteúdo das mensagens e, quando possível, os metadados da comunicação, frequentemente expostos nos cabeçalhos (\textit{headers}) dos protocolos. Esses metadados podem revelar padrões de comunicação e outras informações que não deveriam ser acessíveis a terceiros não autorizados.

Este trabalho tem origem em um projeto desenvolvido em conjunto com a DATACOM, cujo objetivo era viabilizar a comunicação entre aplicações situadas em redes distintas, com o mínimo de infraestrutura intermediária. O projeto explorou técnicas de travessia de NAT (\textit{NAT traversal}) e o tunelamento de tráfego TCP sobre UDP, com base no Chownat. A partir desse desenvolvimento, o presente trabalho aprofunda a integração do QUIC com mecanismos de estabelecimento de conectividade baseados em ICE e considera sua aplicação à comunicação entre \textit{gateways} de sistemas IoT.

A solução proposta consiste em uma infraestrutura de comunicação ponto a ponto entre \textit{gateways} de sistemas IoT situados em redes distintas. A infraestrutura mantém as aplicações baseadas em TCP em conexões locais, transporta seus dados entre os \textit{gateways} sobre o protocolo QUIC \citep{rfc9000} e utiliza o protocolo ICE \citep{rfc8445} para descobrir e selecionar os caminhos de comunicação. A proposta visa aprimorar o seguinte conjunto de atributos de segurança e dependabilidade na transmissão de dados em aplicações IoT:

\begin{itemize}
    \item \textbf{Confidencialidade:} garantia de que os dados transmitidos sejam acessíveis somente a entidades autorizadas.
    \item \textbf{Confiabilidade:} capacidade de um sistema de fornecer corretamente o serviço esperado, de forma contínua, durante um determinado período.
    \item \textbf{Disponibilidade:} prontidão do sistema para fornecer o serviço corretamente quando solicitado.
    \item \textbf{Integridade:} ausência de alterações indevidas no sistema, incluindo modificações não autorizadas nos dados transmitidos.
\end{itemize}

A Figura~\ref{fig:quic_p2p_architecture} apresenta uma visão geral da solução. A arquitetura considerada neste trabalho é composta por dois \textit{peers} e por servidores auxiliares empregados pelos mecanismos de travessia de NAT. Cada \textit{peer} é um \textit{gateway} associado a um ou mais dispositivos IoT ou a uma aplicação consumidora dos dados gerados por esse dispositivo, representados na figura como Dispositivo~1 e Dispositivo~2. Essa aplicação pode ser executada em um sistema de propósito geral ou em outro dispositivo IoT. Os \textit{gateways} executam a infraestrutura proposta e são distintos dos roteadores de acesso, responsáveis pelo encaminhamento dos pacotes e pela tradução de endereços nas redes representadas. As aplicações locais assumem os papéis de cliente e servidor, enquanto os \textit{gateways} viabilizam a comunicação entre elas.

Na figura, a linha tracejada indica o caminho QUIC direto entre os \textit{gateways}, estabelecido por \textit{hole punching}. Os servidores \textit{relay} oferecem o caminho alternativo quando a conexão direta não é possível.

\begin{figure}[!ht]
    \centering
    % Visão geral da arquitetura proposta (Figura fig:quic_p2p_architecture)
% Requer: \usepackage{tikz} e \usetikzlibrary{positioning,calc,shapes.geometric,shapes.symbols}
\resizebox{\linewidth}{!}{%
\begin{tikzpicture}[
  font=\sffamily\small,
  node distance=6mm,
  gw/.style={draw, circle, minimum size=15mm, align=center, inner sep=1pt, fill=blue!8},
  dev/.style={draw, rectangle, minimum width=17mm, minimum height=8mm, align=center, fill=gray!10},
  box/.style={draw, rectangle, minimum height=10mm, minimum width=13mm, align=center, inner sep=2pt},
  srv/.style={draw, ellipse, minimum width=22mm, minimum height=11mm, align=center, inner sep=1pt, fill=orange!10},
  net/.style={draw, cloud, cloud puffs=10, cloud puff arc=110, aspect=1.6, minimum width=17mm, minimum height=10mm, align=center, inner sep=0pt, fill=white},
  link/.style={draw, thick}
]

% Eixo principal
\node[gw] (ga) at (0,0) {\textit{Gateway}\\A};
\node[net] (i1) at (5.9,0) {Internet};
\node[srv] (sig) at (8.6,0) {Servidor de\\sinalização};
\node[net] (i2) at (11.3,0) {Internet};
\node[gw] (gb) at (17.3,0) {\textit{Gateway}\\B};

% Roteadores e NATs centralizados entre Gateway e Internet
% Lado A: Gateway A -> Roteador A -> NAT A -> Internet
\coordinate (midA) at ($(ga.center)!0.5!(i1.center)$);
\node[box, fill=green!8, anchor=east] (r1) at (midA) {Roteador\\A};
\node[box, fill=red!8, anchor=west] (n1) at (r1.east) {NAT\\A};

% Lado B: Internet -> NAT B -> Roteador B -> Gateway B
\coordinate (midB) at ($(i2.center)!0.5!(gb.center)$);
\node[box, fill=red!8, anchor=east] (n2) at (midB) {NAT\\B};
\node[box, fill=green!8, anchor=west] (r2) at (n2.east) {Roteador\\B};

% Dispositivos
\node[dev, below=12mm of ga] (d1) {Dispositivo 1};
\node[dev, below=12mm of gb] (d2) {Dispositivo 2};

% Servidores auxiliares
\node[srv] (stun) at (8.6,2.4) {Servidor\\STUN};
\node[srv] (rl1) at (6.2,-2.6) {Servidor\\\textit{relay} 1};
\node[net] (i3) at (8.6,-2.6) {Internet};
\node[srv] (rl2) at (11.0,-2.6) {Servidor\\\textit{relay} 2};

% Ligações
\draw[link] (ga) -- (r1);
\draw[link] (n1) -- (i1);
\draw[link] (i1) -- (sig);
\draw[link] (sig) -- (i2);
\draw[link] (i2) -- (n2);
\draw[link] (r2) -- (gb);
\draw[link] (ga) -- (d1);
\draw[link] (gb) -- (d2);
\draw[link] (i1.north) |- (stun.west);
\draw[link] (i2.north) |- (stun.east);
\draw[link] (i1.south) |- (rl1.north);
\draw[link] (i2.south) |- (rl2.north);
\draw[link] (rl1) -- (i3);
\draw[link] (i3) -- (rl2);

% Caminho QUIC direto
\draw[thick, dashed, black!70] (ga.north) .. controls (2,4.6) and (15.3,4.6) .. (gb.north)
  node[pos=0.5, above, font=\sffamily\small, align=center] {Comunicação P2P direta com QUIC};

\end{tikzpicture}%
}
    \caption{Visão geral da arquitetura proposta: caminho QUIC direto entre os \textit{gateways} e servidores auxiliares de sinalização, STUN e \textit{relay}.}
    \label{fig:quic_p2p_architecture}
\end{figure}

\section{Motivação}

Em sistemas IoT organizados por meio de \textit{gateways}, esses nós assumem a responsabilidade pela comunicação externa das aplicações associadas aos dispositivos locais. Quando situados em redes privadas distintas, os \textit{gateways} precisam de mecanismos que permitam estabelecer conectividade através de NATs. Nesse cenário, a comunicação ponto a ponto pode reduzir a dependência de servidores intermediários para o transporte contínuo dos dados, mantendo serviços auxiliares para a descoberta e o estabelecimento das conexões.

A comunicação direta entre pares (\textit{peers}) situados em redes privadas distintas depende de técnicas de travessia de NAT \citep{rfc5128}, isto é, técnicas que permitem a comunicação através de tradutores de endereços de rede (NAT, do inglês \textit{Network Address Translator}). Os protocolos STUN (\textit{Session Traversal Utilities for NAT}), TURN (\textit{Traversal Using Relays around NAT}) e ICE (\textit{Interactive Connectivity Establishment}), padronizados pelo IETF, implementam essas técnicas. O WebRTC (\textit{Web Real-Time Communication}) \citep{WebRTCForTheCurious}, padrão aberto para comunicação em tempo real entre navegadores, é o exemplo mais difundido do uso conjunto desses protocolos. Por isso, este trabalho adota o WebRTC como referência arquitetural para a organização da sinalização e do estabelecimento de conectividade. A solução proposta, entretanto, não utiliza a pilha WebRTC: ela combina diretamente o ICE, o STUN e o TURN com o QUIC na comunicação entre \textit{gateways}.

O QUIC \citep{rfc9000}, por sua vez, é um protocolo de transporte que opera sobre o UDP (\textit{User Datagram Protocol}) e foi desenvolvido para superar limitações da comunicação baseada em TCP (\textit{Transmission Control Protocol}). Sua concepção está associada à redução da latência e à integração de mecanismos de segurança, inicialmente motivadas pelas necessidades das conexões HTTP (\textit{Hypertext Transfer Protocol}) na Web \citep{marx2020quic}. O QUIC implementa sobre UDP mecanismos de entrega confiável e controle de congestionamento, além de integrar mecanismos de proteção criptográfica ao estabelecimento da conexão. Sua implementação no espaço de usuário também facilita a evolução do protocolo sem depender de alterações no núcleo (\textit{kernel}) dos sistemas operacionais. Essas características motivam a investigação do QUIC como transporte externo entre os \textit{gateways}. A proposta mantém a interface TCP local das aplicações existentes e transporta seus dados pelo caminho estabelecido com ICE, buscando facilitar a integração sem exigir mudanças no protocolo utilizado por essas aplicações.

\section{Objetivos}

Este trabalho tem por objetivo desenvolver e avaliar uma infraestrutura de comunicação ponto a ponto entre \textit{gateways} de sistemas IoT situados em redes distintas, capaz de integrar aplicações legadas baseadas em TCP a um transporte externo sobre QUIC. A solução utiliza o ICE para descobrir e selecionar caminhos de comunicação entre os \textit{gateways}. A proposta mantém as conexões TCP das aplicações na interface local de \textit{loopback} e transporta seus dados via QUIC até o \textit{gateway} remoto, de modo a preservar a interface de comunicação das aplicações e reduzir a necessidade de modificações em sua implementação.

Os atributos de confidencialidade, integridade, confiabilidade e disponibilidade apresentados anteriormente orientam a discussão dos mecanismos utilizados e das limitações da solução, considerando as configurações e os cenários efetivamente avaliados. Para alcançar o objetivo geral, são definidos os seguintes objetivos específicos:

\begin{itemize}
\item Desenvolver uma interface modular de comunicação que permita integrar diferentes protocolos de transporte e mecanismos de estabelecimento de conectividade, tomando como ponto de partida a arquitetura do Chownat.
\item Implementar o transporte de dados das aplicações TCP locais sobre QUIC, utilizando a biblioteca \texttt{OpenSSL}, de modo a preservar a interface de comunicação dessas aplicações nos \textit{gateways}.
\item Integrar o protocolo ICE por meio da biblioteca \texttt{libnice}, permitindo a descoberta e a seleção de caminhos de comunicação entre os \textit{gateways}, incluindo caminhos diretos e caminhos intermediados por servidores TURN.
\item Verificar o funcionamento da infraestrutura em topologias de rede controladas e sob diferentes configurações de NAT, considerando o estabelecimento de conectividade por \textit{hole punching} UDP ou por retransmissão via TURN, a transferência de dados e o encerramento das conexões.
\item Comparar conceitualmente a abordagem proposta com a comunicação mediada por servidor e com a comunicação P2P baseada em TCP \textit{hole punching}.
\end{itemize}

O escopo deste trabalho concentra-se na comunicação externa entre \textit{gateways} de uma arquitetura IoT. Nos experimentos, esses \textit{gateways} são representados por \textit{hosts} de propósito geral que executam a infraestrutura proposta e aplicações locais de geração e recepção de mensagens com dados sintéticos. O trabalho não avalia a aquisição de dados por sensores, a comunicação entre dispositivos IoT e \textit{gateways}, a execução das bibliotecas em microcontroladores nem a eficiência energética desses dispositivos.

\section{Organização do Documento}

Este documento está organizado da seguinte forma:

\begin{itemize}
\item O Capítulo~\ref{chap:background} apresenta os fundamentos de comunicação necessários à compreensão da solução, incluindo os modelos de arquitetura de rede, o papel dos \textit{gateways}, o funcionamento do NAT, as técnicas de travessia de NAT e os protocolos STUN, TURN e ICE. Também apresenta o Chownat e a relação da proposta com a arquitetura do WebRTC.
\item O Capítulo~\ref{chap:quic} descreve o protocolo QUIC, abordando sua estrutura, os mecanismos de estabelecimento e manutenção das conexões, o transporte de dados e sua utilização em conjunto com técnicas de travessia de NAT.
\item O Capítulo~\ref{chap:implementation} detalha a implementação da infraestrutura proposta: a extensão do Chownat, a interface modular de comunicação, o tunelamento das aplicações TCP locais e a integração das bibliotecas \texttt{OpenSSL} e \texttt{libnice}. O capítulo também descreve as topologias e os ambientes experimentais.
\item O Capítulo~\ref{chap:results} apresenta a metodologia de avaliação, os resultados dos testes funcionais e a análise conceitual das abordagens de comunicação, e discute o comportamento da infraestrutura nos cenários considerados e as limitações da avaliação experimental.
\item O Capítulo~\ref{chap:conclusion} apresenta as conclusões do trabalho, relaciona os resultados aos objetivos estabelecidos e indica possibilidades de continuidade.
\end{itemize}
